% ------------------------------------------------------------
% Appendix: Deriving the (negative) ELBO loss for the SeqGPLVM
% ------------------------------------------------------------

\section{SeqGPLVM loss function}
\label{sec:seqgplvm_loss}
In this section, we derive the ELBO (evidence lower bound) for the sequential Gaussian process latent variable 
model (SeqGPLVM) introduced in Section~\ref{sec:seqgplvm}. 

\subsection{Notation}
We use bold uppercase letters to denote matrices, bold lowercase
letters to denote vectors, and plain letters to denote scalars. Let $\mathbf{X} \in \mathbb{R}^{N\times T \times D}$
denote the covariate data 
where \(N\) is the number of individuals and 
\(D\) is the number of covariates and \(T\) is the number of time points, and $\mathbf{A} \in \{0,1\}^{N \times T}$ the observed outcomes (treatments). Here, 
we assume scalar outcomes for simplicity, but the 
extension to multi-dimensional outcomes is straightforward. Also, we assume 
a binary outcome for concreteness; 
however, this restriction is not required 
for the theoretical results. We index individuals 
by \(i \in \{1,\dots,N\}\) and time by \(t \in \{1,\dots,T\}\).
Bold lowercase symbols (e.g., $\mathbf{a}_t$) denote vectors collecting
observations for all $N$ units at time $t$, while plain symbols (e.g., $x_{i,d,t}$
and $a_{i,t}$) denote scalar observations for unit $i$ at time $t$. We use $\boldsymbol{\lambda}$ to denote 
the vector of \textit{all} latent variables which we want to calculate the 
posterior, including the \emph{time-invariant} latent confounder $\mathbf{Z} \in \mathbb{R}^{N \times P}$, where 
\(P\) is the latent dimensionality,
and $\boldsymbol{\gamma}$ to denote the \emph{generative} hyperparameters and $\boldsymbol{\psi}$ 
to denote the \emph{variational} hyperparameters. Generative hyperparameters are those that define 
how the data is generated and they are in prior and likelihood, 
e.g., the kernel lengthscale and variance in GPs, 
while variational hyperparameters are those that define the 
variational family, e.g., the mean and variance in Gaussian 
variational distributions. 

\subsection{Variational inference background}
\label{sec:variational_inference_background}
Our inferential target is the joint posterior distribution over latent variables given the data and hyperparameters:
\(
p_{\boldsymbol{\gamma}}(\boldsymbol{\lambda}\mid \mathbf{X},\mathbf{A})
\). Bayes' rule gives 

\begin{equation}
p_{\boldsymbol{\gamma}}(\boldsymbol{\lambda}\mid \mathbf{X},\mathbf{A})
=
\frac{p_{\boldsymbol{\gamma}}(\mathbf{A}\mid \mathbf{X},\boldsymbol{\lambda})\,p_{\boldsymbol{\gamma}}(\boldsymbol{\lambda}\mid \mathbf{X})}
     {p_{\boldsymbol{\gamma}}(\mathbf{A}\mid \mathbf{X})},
\label{eq:bayes-posterior}
\end{equation}
where the denominator
\(
p_{\boldsymbol{\gamma}}(\mathbf{A}\mid \mathbf{X})
\)
(the \emph{marginal likelihood} / \emph{evidence}) is the normalizing constant. Note that 
we show the parameterization both likehood and prior with $\boldsymbol{\gamma}$ to make the 
notation consicer, but in practice, they can have different hyperparameters.
Rewriting the logarithm of evidence as an 
integral over latent variables,  
\begin{equation}
\log p_{\boldsymbol{\gamma}}(\mathbf{A}\mid \mathbf{X}) = \log \int p_{\boldsymbol{\gamma}}(\mathbf{A},\boldsymbol{\lambda}\mid \mathbf{X})\,\mathrm{d}\boldsymbol{\lambda}.
\end{equation}
In latent-variable models this integral is typically intractable which means the 
posterior cannot be computed exactly; introducing a tractable variational density \(q_{\boldsymbol{\psi}}(\boldsymbol{\lambda}\mid \mathbf{X})\) 
yields a lower bound on \(\log p_{\boldsymbol{\gamma}}(\mathbf{A}\mid\mathbf{X})\). 
The gap between this bound and the true log-evidence is a KL divergence to the true posterior, 
so maximizing the bound pushes \(q_{\boldsymbol{\psi}}(\boldsymbol{\lambda}\mid \mathbf{X})\) 
towards \(p_{\boldsymbol{\gamma}}(\boldsymbol{\lambda}\mid \mathbf{X},\mathbf{A})\). To see this, 
the marginal likelihood can be written as
\begin{align}
\log p_{\boldsymbol{\gamma}}(\mathbf{A}\mid \mathbf{X})
&= \log \int p_{\boldsymbol{\gamma}}(\mathbf{A},\boldsymbol{\lambda}\mid \mathbf{X})\,\mathrm{d}\boldsymbol{\lambda}
\nonumber\\
&= \log \int q_{\boldsymbol{\psi}}(\boldsymbol{\lambda}\mid \mathbf{X})\,
\frac{p_{\boldsymbol{\gamma}}(\mathbf{A},\boldsymbol{\lambda}\mid \mathbf{X})}{q_{\boldsymbol{\psi}}(\boldsymbol{\lambda}\mid \mathbf{X})}\,
\mathrm{d}\boldsymbol{\lambda}
\nonumber\\
&\ge \int q_{\boldsymbol{\psi}}(\boldsymbol{\lambda}\mid \mathbf{X})
\log\frac{p_{\boldsymbol{\gamma}}(\mathbf{A},\boldsymbol{\lambda}\mid \mathbf{X})}{q_{\boldsymbol{\psi}}(\boldsymbol{\lambda}\mid \mathbf{X})}\,
\mathrm{d}\boldsymbol{\lambda}
\;=:\; \mathcal{L}(\boldsymbol{\psi},\boldsymbol{\gamma}),
\label{eq:elbo-def}
\end{align}
where the inequality is Jensen's inequality. $\mathcal{L}(\boldsymbol{\psi},\boldsymbol{\gamma})$ is called the evidence lower bound (ELBO).
Using Bayes' rule,
\(
p_{\boldsymbol{\gamma}}(\boldsymbol{\lambda}\mid \mathbf{X},\mathbf{A})
=
p_{\boldsymbol{\gamma}}(\mathbf{A},\boldsymbol{\lambda}\mid \mathbf{X})/p_{\boldsymbol{\gamma}}(\mathbf{A}\mid \mathbf{X})
\),
one obtains the standard identity
\begin{align} \mathrm{KL}\!\left(q_{\boldsymbol{\psi}}(\boldsymbol{\lambda}\mid \mathbf{X})\,\|\,p_{\boldsymbol{\gamma}}(\boldsymbol{\lambda}\mid \mathbf{X},\mathbf{A})\right) 
    &= \int q_{\boldsymbol{\psi}}(\boldsymbol{\lambda}\mid \mathbf{X}) \log\frac{q_{\boldsymbol{\psi}}(\boldsymbol{\lambda}\mid \mathbf{X})}{p_{\boldsymbol{\gamma}}(\boldsymbol{\lambda}\mid \mathbf{X},\mathbf{A})}\, \mathrm{d}\boldsymbol{\lambda} \nonumber\\ 
    &= \int q_{\boldsymbol{\psi}}(\boldsymbol{\lambda}\mid \mathbf{X}) \left [ \log q_{\boldsymbol{\psi}}(\boldsymbol{\lambda}\mid \mathbf{X}) - \log {p_{\boldsymbol{\gamma}}(\boldsymbol{\lambda}\mid \mathbf{X},\mathbf{A})}\right ] \, \mathrm{d}\boldsymbol{\lambda} \nonumber\\ 
    &= \int q_{\boldsymbol{\psi}}(\boldsymbol{\lambda}\mid \mathbf{X}) \left [ \log q_{\boldsymbol{\psi}}(\boldsymbol{\lambda}\mid \mathbf{X}) - \log\frac{p_{\boldsymbol{\gamma}}(\mathbf{A},\boldsymbol{\lambda}\mid \mathbf{X})}{p_{\boldsymbol{\gamma}}(\mathbf{A}\mid \mathbf{X})}\right ] \, \mathrm{d}\boldsymbol{\lambda} \nonumber\\ 
    &= \log p_{\boldsymbol{\gamma}}(\mathbf{A}\mid \mathbf{X}) -\int q_{\boldsymbol{\psi}}(\boldsymbol{\lambda}\mid \mathbf{X}) \log\frac{p_{\boldsymbol{\gamma}}(\mathbf{A},\boldsymbol{\lambda}\mid \mathbf{X})}{q_{\boldsymbol{\psi}}(\boldsymbol{\lambda}\mid \mathbf{X})}\, \mathrm{d}\boldsymbol{\lambda} \nonumber\\ 
    &= \log p_{\boldsymbol{\gamma}}(\mathbf{A}\mid \mathbf{X}) - \mathcal{L}(\boldsymbol{\psi},\boldsymbol{\gamma}). 
    \label{eq:logp-elbo-kl} 
\end{align}
Equivalently,
\begin{equation}
\log p_{\boldsymbol{\gamma}}(\mathbf{A}\mid \mathbf{X})
=
\mathcal{L}(\boldsymbol{\psi},\boldsymbol{\gamma})
+
\mathrm{KL}\!\left(q_{\boldsymbol{\psi}}(\boldsymbol{\lambda})\,\|\,p_{\boldsymbol{\gamma}}(\boldsymbol{\lambda}\mid \mathbf{X},\mathbf{A})\right),
\label{eq:logp-decomp}
\end{equation}
so maximizing \(\mathcal{L}(\boldsymbol{\psi},\boldsymbol{\gamma})\) w.r.t.\ \(q\) minimizes the KL divergence to the true posterior (and equality holds iff \(q_{\boldsymbol{\psi}}(\boldsymbol{\lambda}) = p_{\boldsymbol{\gamma}}(\boldsymbol{\lambda}\mid \mathbf{X},\mathbf{A})\)). 
This is true since the lhs of \eqref{eq:logp-decomp} does not depend on \(q\). 
Thus, we can use \(\mathcal{L}(\boldsymbol{\psi},\boldsymbol{\gamma})\) as a surrogate objective for approximate inference. 
We can rewrite \(\mathcal{L}(\boldsymbol{\psi},\boldsymbol{\gamma})\) as 
\begin{align}
\mathcal{L}(\boldsymbol{\psi},\boldsymbol{\gamma})
&= \int q_{\boldsymbol{\psi}}(\boldsymbol{\lambda})\,\log\!\left(\frac{p_{\boldsymbol{\gamma}}(\mathbf{A},\boldsymbol{\lambda} \mid \mathbf{X})}{q_{\boldsymbol{\psi}}(\boldsymbol{\lambda})}\right)\, d\boldsymbol{\lambda} \\
&= \int q_{\boldsymbol{\psi}}(\boldsymbol{\lambda})\,\log p_{\boldsymbol{\gamma}}(\mathbf{A},\boldsymbol{\lambda} \mid \mathbf{X})\, d\boldsymbol{\lambda}
   \;-\; \int q_{\boldsymbol{\psi}}(\boldsymbol{\lambda})\,\log q_{\boldsymbol{\psi}}(\boldsymbol{\lambda})\, d\boldsymbol{\lambda} \\
&= \int q_{\boldsymbol{\psi}}(\boldsymbol{\lambda})\,\log\!\Big( p_{\boldsymbol{\gamma}}(\mathbf{A}\mid \mathbf{X},\boldsymbol{\lambda})\,p_{\boldsymbol{\gamma}}(\boldsymbol{\lambda}\mid \mathbf{X}) \Big)\, d\boldsymbol{\lambda}
   \;-\; \int q_{\boldsymbol{\psi}}(\boldsymbol{\lambda})\,\log q_{\boldsymbol{\psi}}(\boldsymbol{\lambda})\, d\boldsymbol{\lambda} \\
&= \int q_{\boldsymbol{\psi}}(\boldsymbol{\lambda})\,\log p_{\boldsymbol{\gamma}}(\mathbf{A}\mid \mathbf{X},\boldsymbol{\lambda})\, d\boldsymbol{\lambda}
   \;-\; \int q_{\boldsymbol{\psi}}(\boldsymbol{\lambda})\,\log\!\left(\frac{q_{\boldsymbol{\psi}}(\boldsymbol{\lambda})}{p_{\boldsymbol{\gamma}}(\boldsymbol{\lambda}\mid \mathbf{X})}\right)\, d\boldsymbol{\lambda} \\
&= \int q_{\boldsymbol{\psi}}(\boldsymbol{\lambda})\,\log p_{\boldsymbol{\gamma}}(\mathbf{A}\mid \mathbf{X},\boldsymbol{\lambda})\, d\boldsymbol{\lambda}
    \;-\; \mathrm{KL}\!\left(q_{\boldsymbol{\psi}}(\boldsymbol{\lambda})\,\|\,p_{\boldsymbol{\gamma}}(\boldsymbol{\lambda}\mid \mathbf{X})\right).
    \label{eq:elbo-rewrite}
\end{align}
Thus to setup variational inference, we need to specify the model likelihood \(p_{\boldsymbol{\gamma}}(\mathbf{A}\mid \mathbf{X},\boldsymbol{\lambda})\), 
the prior over latents \(p_{\boldsymbol{\gamma}}(\boldsymbol{\lambda}\mid \mathbf{X})\), 
and the variational family \(q_{\boldsymbol{\psi}}(\boldsymbol{\lambda})\).

\subsection{Variational Inference for Bayesian Gaussian process latent variable models}
In this section, we briefly review variational inference in \emph{Bayesian} Gaussian 
process latent variable models (GPLVMs) \cite{titsiasBayesianGaussianProcess2010}. The word Bayesian 
here is to distinguish it from the original GPLVM \cite{lawrenceGaussianProcessLatent2003} 
where the latent variables \( \mathbf{Z} \) are point estimated rather than 
inferred in a fully Bayesian manner. From now on, by GPLVMs we mean Bayesian GPLVMs. 
The idea behind GPLVMs is to place a Gaussian processes (GPs) as a 
prior in Eq.~\ref{eq:elbo-rewrite}. To show the machinary 
of GPLVMs, first forget about the time index \(t\) or 
equivalently think of slice the data at a fixed time \(t' \). We 
define \(\widetilde{\boldsymbol{X}}_{t'} \in \mathbb{R}^{N \times (D+p)}\) 
as the augmented input defined 
on the joint space of observed covariates and latent variables: 
\begin{equation}
    \widetilde{\boldsymbol{x}}_{i,t'} = (\boldsymbol{x}_{i,t'}, \mathbf{z}_i), \quad i=1,\dots,N.
\end{equation}
Then, we define a GP prior over the latent function \(f_{t'}(\cdot)\) that maps 
these augmented inputs to the observed outcomes:
\begin{equation}
    f_{t'}(\cdot) \sim \mathcal{GP}\left(m_{t'}(\cdot), k_{t'}(\cdot, \cdot; \boldsymbol{\gamma}_{t'})\right),
\end{equation}
where \(m_{t'}(\cdot)\) is the mean function, \(k_{t'}(\cdot, \cdot; \boldsymbol{\gamma}_{t'})\) is the kernel function 
parameterized by hyperparameters \(\boldsymbol{\gamma}_{t'}\). We then perform inference 
over the finite vector of function values,
\begin{equation}
\mathbf{f}_{t'} = f_{t'}(\widetilde{\mathbf{X}}_{t'}) \in \mathbb{R}^{N}
\end{equation}
Hence, the latent variables in GPLVMs are 
\( \boldsymbol{\lambda}_{t'} = (\mathbf{f}_{t'}, \mathbf{Z}) \) and 
the posterior we are interested in is \(p_{\boldsymbol{\gamma}}(\mathbf{f}_{t'}, \mathbf{Z} \mid \mathbf{X}_{t'}, \mathbf{Y}_{t'})\). 
As we described in \ref{sec:variational_inference_background}, to perform variational inference, 
we need to specify the prior over latents \(p_{\boldsymbol{\gamma}}(\boldsymbol{\lambda}_{t'}\mid \mathbf{X}_{t'})\). In GPLVMs,
we can decompose the prior as follows: 
\begin{align}
    p_{\boldsymbol{\gamma}}(\boldsymbol{\lambda}_{t'}\mid \mathbf{X}_{t'}) 
    &= p_{\boldsymbol{\gamma}}(\mathbf{f}_{t'}, \mathbf{Z} \mid \mathbf{X}_{t'}) \\
    &= p_{\boldsymbol{\gamma}}(\mathbf{f}_{t'} \mid \mathbf{X}_{t'}, \mathbf{Z}) p_{\boldsymbol{\gamma}}(\mathbf{Z}\mid \mathbf{X}_{t'}),
\end{align}
where the first term is the GP prior conditional on the inputs. 
In vanilla GPLVMs, 
we assume that the latent variables are \emph{a priori} 
independent of the observed covariates, and that they are i.i.d.\ 
with a standard normal prior, e.g., \(\mathcal{N}(\mathbf{0}, \mathbf{I}_P)\).
\begin{align}
p_{\boldsymbol{\gamma}}(\mathbf{Z}\mid \mathbf{X}_{t'})
&= p(\mathbf{Z})\\
&= \prod_{i=1}^N \mathcal{N}(\mathbf{z}_i \mid \mathbf{0}, \mathbf{I}_P).
\label{eq:gplvm_latent_prior}
\end{align}

Next, we need to specify the variational family $q_{\boldsymbol{\psi}}(\boldsymbol{\lambda}_{t'}\mid \mathbf{X}_{t'})$. 
Following the standard construction in variational inference for Gaussian process models, 
we define a partially collapsed variational family in which only the latent variables $\mathbf{Z}$ 
are approximated variationally, while the conditional distribution of the GP function values is 
kept exact. Shortly we will see the advantage of this parameterization when we derive the ELBO. Specifically, we set
\begin{align}
    q_{\boldsymbol{\psi},\boldsymbol{\gamma}}(\boldsymbol{\lambda}_{t'}\mid \mathbf{X}_{t'}) 
    &= q_{\boldsymbol{\psi},\boldsymbol{\gamma}}(\mathbf{f}_{t'}, \mathbf{Z}\mid \mathbf{X}_{t'}) \\
    &= p_{\boldsymbol{\gamma}}(\mathbf{f}_{t'} \mid \mathbf{Z}, \mathbf{X}_{t'}) q_{\boldsymbol{\psi}}(\mathbf{Z}\mid \mathbf{X}_{t'}) \\
    &= p_{\boldsymbol{\gamma}}(\mathbf{f}_{t'} \mid \mathbf{Z}, \mathbf{X}_{t'}) \prod_{i=1}^N q_{\boldsymbol{\psi}}(\mathbf{z}_i),
    \label{eq:gplvm_var_posterior}
\end{align}  
 
Note that the variational distribution is indexed by both 
\(\boldsymbol{\psi}\) and \(\boldsymbol{\gamma}\) because of this choice unlike the usual 
variational inference setting. See section 4.1 \cite{leibfriedTutorialSparseGaussian2022} for more details.The factorization in Eq.~\ref{eq:gplvm_var_posterior} is known as 
the mean field approximation. It assumes that the 
latent variables are independent across data points in the variational posterior. We also assume that 
the variational distribution over the latent variables is not conditioned on the observed covariates, i.e.,
\(q_{\boldsymbol{\psi}}(\mathbf{Z}\mid \mathbf{X}_{t'}) = q_{\boldsymbol{\psi}}(\mathbf{Z})\). This is called 
the none-amortized variational inference. Moreover, we assume 
a fully factorized Gaussian variational distribution for each latent variable:
\begin{equation}
    q_{\boldsymbol{\psi}}(\mathbf{z}_i) = \prod_{j=1}^P \mathcal{N}(z_{ij} \mid \mu_{ij}, \sigma_{ij}^2).
\end{equation}

substituting the prior and variational family into the ELBO in Eq.~\ref{eq:elbo-rewrite}, we obtain
\begin{align}
\mathcal{L}_{t'}(\boldsymbol{\psi},\boldsymbol{\gamma})
&= \int q_{\boldsymbol{\psi},\boldsymbol{\gamma}}(\boldsymbol{\lambda}_{t'})\,\log p_{\boldsymbol{\gamma}}(\mathbf{A}_{t'}\mid \mathbf{X}_{t'},\boldsymbol{\lambda}_{t'})\, d\boldsymbol{\lambda}_{t'}
    \;-\; \mathrm{KL}\!\left(q_{\boldsymbol{\psi},\boldsymbol{\gamma}}(\boldsymbol{\lambda}_{t'})\,\|\,p_{\boldsymbol{\gamma}}(\boldsymbol{\lambda}_{t'}\mid \mathbf{X}_{t'})\right) \\
&= \int q_{\boldsymbol{\psi}}(\mathbf{Z})\,p_{\boldsymbol{\gamma}}(\mathbf{f}_{t'} \mid \mathbf{Z})\,\log p_{\boldsymbol{\gamma}}(\mathbf{A}_{t'}\mid \mathbf{X}_{t'},\mathbf{f}_{t'})\, d\mathbf{f}_{t'} d\mathbf{Z}
    \;-\; \mathrm{KL}\!\left(q_{\boldsymbol{\psi}}(\mathbf{Z})\,\|\,p_{\boldsymbol{\gamma}}(\mathbf{Z}\mid \mathbf{X}_{t'})\right) \\
&= \sum_{i=1}^N \int q_{\boldsymbol{\psi}}(\mathbf{Z})\,p_{\boldsymbol{\gamma}}(\mathbf{f}_{t'} \mid \mathbf{Z})\,\log p_{\boldsymbol{\gamma}}(a_{i,t'}\mid,f_{i,t'})\,d\mathbf{Z}
    \;-\; \mathrm{KL}\!\left(q_{\boldsymbol{\psi}}(\mathbf{Z})\,\|\,p_{\boldsymbol{\gamma}}(\mathbf{Z}\mid \mathbf{X}_{t'})\right) \\
&= \mathbb{E}_{q_{\boldsymbol{\psi}}(\mathbf{Z})\,p_{\boldsymbol{\gamma}}(\mathbf{f}_{t'} \mid \mathbf{Z})}
\left[\log p_{\boldsymbol{\gamma}}(\mathbf{A}_{t'}\mid \mathbf{X}_{t'},\mathbf{f}_{t'})\right]
    \;-\; \mathrm{KL}\!\left(q_{\boldsymbol{\psi}}(\mathbf{Z})\,\|\,p_{\boldsymbol{\gamma}}(\mathbf{Z}\mid \mathbf{X}_{t'})\right).
\label{eq:gplvm_elbo}
\end{align}

\subsection{Sparse GPLVMs}
The ELBO in Eq.~\ref{eq:gplvm_elbo} is still computationally expensive to evaluate
because it involves the full GP prior \(p_{\boldsymbol{\gamma}}(\mathbf{f}_{\widetilde{\mathbf{X}}} \mid \mathbf{Z})\), which 
scales cubically with the number of data points \(N\). To address this issue, we can
\subsection{SeqGPLVM: inducing variables and the resulting ELBO}

\paragraph{Inducing variables and GP conditionals.}
For each time step \(t\in\{1,\dots,T\}\), let \(\mathbf{f}_t\in\mathbb{R}^{N}\) denote latent GP function values at the \(N\) observed inputs and \(\mathbf{u}_t\in\mathbb{R}^{M}\) the inducing variables at inducing inputs \(\mathbf{M}_t\) (collecting \(M\) inducing locations). With kernel hyperparameters \(\boldsymbol{\psi}_t\), define kernel matrices
\(
\mathbf{K}^{(t)}_{NN}, \mathbf{K}^{(t)}_{NM}, \mathbf{K}^{(t)}_{MN}, \mathbf{K}^{(t)}_{MM}
\)
computed from the (time-\(t\)) GP inputs (in the GPLVM these inputs may involve \(\mathbf{X}\) and the latent \(\mathbf{z}\); we suppress this dependence in the conditioning to keep notation light).
The sparse GP prior uses
\begin{align}
p(\mathbf{u}_t\mid \boldsymbol{\psi}_t) &= \mathcal{N}\!\left(\mathbf{0},\mathbf{K}^{(t)}_{MM}\right),
\\
p(\mathbf{f}_t\mid \mathbf{u}_t,\boldsymbol{\psi}_t)
&=\mathcal{N}\!\left(
\mathbf{K}^{(t)}_{NM}(\mathbf{K}^{(t)}_{MM})^{-1}\mathbf{u}_t,\;
\mathbf{K}^{(t)}_{NN}-\mathbf{K}^{(t)}_{NM}(\mathbf{K}^{(t)}_{MM})^{-1}\mathbf{K}^{(t)}_{MN}
\right).
\label{eq:gp-conditional}
\end{align}

\paragraph{Latent variables and variational family.}
Let \(\mathbf{z}\in\mathbb{R}^{d}\) denote the (global) latent variable(s) with prior
\begin{equation}
p(\mathbf{z})=\mathcal{N}(\mathbf{0},\mathbf{I}).
\label{eq:z-prior}
\end{equation}
Collect all latents as
\(
\boldsymbol{\lambda} := (\mathbf{z},\mathbf{f}_{1:T},\mathbf{u}_{1:T})
\),
where \(\mathbf{f}_{1:T}:=(\mathbf{f}_1,\dots,\mathbf{f}_T)\) and similarly for \(\mathbf{u}_{1:T}\).
We use the standard sparse variational GP factorization
\begin{equation}
q(\boldsymbol{\lambda})
=
q(\mathbf{z})\prod_{t=1}^{T} q(\mathbf{u}_t)\,p(\mathbf{f}_t\mid \mathbf{u}_t,\boldsymbol{\psi}_t),
\qquad
q(\mathbf{z})=\mathcal{N}(\mathbf{m}_z,\mathbf{S}_z),
\quad
q(\mathbf{u}_t)=\mathcal{N}(\mathbf{m}_{u_t},\mathbf{S}_{u_t}).
\label{eq:q-factorization}
\end{equation}

\paragraph{Sequential likelihood factorization.}
Assume a time-factorized observation model (conditional on the latent state):
\begin{equation}
p(\mathbf{Y}\mid \mathbf{X},\mathbf{z},\mathbf{f}_{1:T},\boldsymbol{\psi})
=
\prod_{t=1}^{T}
p\!\left(\mathbf{y}_t \mid \mathbf{X}_t,\mathbf{Y}_{1:t-1},\mathbf{z},\mathbf{f}_t,\boldsymbol{\psi}_t\right),
\label{eq:seq-likelihood}
\end{equation}
where \(\mathbf{y}_t\in\mathbb{R}^{N}\) is the vector of observations at time \(t\) and \(\mathbf{Y}_{1:t-1}\) collects past observations.

\paragraph{ELBO for the SeqGPLVM.}
Plugging \eqref{eq:q-factorization} and \eqref{eq:seq-likelihood} into the generic ELBO \eqref{eq:elbo-def} yields
\begin{align}
\mathcal{L}(q)
&=
\mathbb{E}_{q(\boldsymbol{\lambda})}
\left[
\log p(\mathbf{Y}\mid \mathbf{X},\boldsymbol{\lambda},\boldsymbol{\psi})
\right]
-
\mathrm{KL}\!\left(q(\boldsymbol{\lambda})\,\|\,p(\boldsymbol{\lambda}\mid \mathbf{X},\boldsymbol{\psi})\right)
\nonumber\\
&=
\sum_{t=1}^{T}
\mathbb{E}_{q(\mathbf{z})\,q(\mathbf{u}_t)\,p(\mathbf{f}_t\mid \mathbf{u}_t,\boldsymbol{\psi}_t)}
\left[
\log p\!\left(\mathbf{y}_t \mid \mathbf{X}_t,\mathbf{Y}_{1:t-1},\mathbf{z},\mathbf{f}_t,\boldsymbol{\psi}_t\right)
\right]
\nonumber\\
&\quad
-\sum_{t=1}^{T}\mathrm{KL}\!\left(q(\mathbf{u}_t)\,\|\,p(\mathbf{u}_t\mid \boldsymbol{\psi}_t)\right)
-\mathrm{KL}\!\left(q(\mathbf{z})\,\|\,p(\mathbf{z})\right).
\label{eq:seqgplvm-elbo}
\end{align}
The corresponding training loss is the negative ELBO,
\begin{equation}
\mathcal{J}_{\text{VI}}(\boldsymbol{\psi})
:=
-\mathcal{L}(q),
\label{eq:neg-elbo-loss}
\end{equation}
optionally optimized jointly over the variational parameters \(\{(\mathbf{m}_z,\mathbf{S}_z),(\mathbf{m}_{u_t},\mathbf{S}_{u_t})\}_{t=1}^T\) and hyperparameters \(\boldsymbol{\psi}\).

\paragraph{Optional: hyperprior regularization.}
If a hyperprior \(p(\boldsymbol{\psi})\) is used, then (up to an additive constant) one targets
\begin{equation}
\log p(\boldsymbol{\psi}\mid \mathbf{X},\mathbf{Y})
\propto
\log p(\mathbf{Y}\mid \mathbf{X},\boldsymbol{\psi}) + \log p(\boldsymbol{\psi})
\;\approx\;
\mathcal{L}(q) + \log p(\boldsymbol{\psi}),
\label{eq:elbo-plus-hyperprior}
\end{equation}
and in the time-separable case \( \log p(\boldsymbol{\psi}) = \sum_{t=1}^{T}\log p(\boldsymbol{\psi}_t)\).

\begin{landscape}
\begin{figure}[p]
  \centering
  \includegraphics[width=\paperwidth,trim=0.5cm 0.5cm 0.5cm 0.5cm,
  clip]{figures/build/fit_seqgplvm_dag.pdf}
  \caption{Inference dag for the SeqGPLVM.}
  \label{fig:seqdecon-dag}
\end{figure}
\end{landscape}



